\begin{abstract}
Comparing two models on a fixed benchmark can end with practical superiority or equivalence, but certifying that decision may require nearly every item. We characterize this cost through distance to the practical decision boundary and the changes still possible in unseen scores. For any uniformly \(\alpha\)-correct adaptive procedure, a coupling argument bounds the probability of leaving a label-changing set unqueried. Its edit-capacity consequence forces at least \((1-2\alpha)M\) expected queries when each of \(M\) coordinates can individually change the correct label. A standard concentration bound gives a sufficient early-stopping regime farther from the boundary; the two bounds describe different scales and are not matching. Controlled replays show that the near-census region widens with paired disagreement, while constant populations remain costly at the boundary despite zero variance. Replays from seven public response matrices place these findings in context: 700 random pairs exhibit heterogeneous costs, whereas 140 retrospectively selected margin-nearest pairs require 97.7--99.7\% of items on average across matrices with direct betting. These stress-test costs demonstrate difficult available comparisons, not their prevalence or the optimality of the evaluator. The resulting account separates an unavoidable obstruction at fine boundary resolution from the broader, method-dependent cost observed in benchmarks.
\end{abstract}
